
Eight conformance profiles are defined in \texttt{profiles/}, including five hierarchical profiles and three regulatory overlays (\texttt{country-us-fda.md}, \texttt{state-us-ca.md}, \texttt{state-us-ny.md}).

\subsection{Safety Modules for Physical AI}

The system provides eight safety modules in \texttt{safety/} specifically designed for robotic systems operating in clinical environments:

\begin{table}[H]
\centering
\caption{Safety modules for Physical AI systems.}
\label{tab:safety-modules}
\footnotesize
\begin{tabularx}{\textwidth}{@{}L{3.8cm}L{8.2cm}@{}}
\toprule
\textbf{Module} & \textbf{Purpose} \\
\midrule
\texttt{estop.py} & Emergency stop signal propagation, abort workflow, post-abort evidence capture, recovery with re-authorization \\
\texttt{procedure\_state.py} & Procedure state machine enforcement (IDLE, PREPARING, ACTIVE, PAUSED, COMPLETING, ABORTED) \\
\texttt{robot\_registry.py} & Robot registration, capability validation, firmware version tracking \\
\texttt{task\_validator.py} & Task order validation against safety constraints and robot capabilities \\
\texttt{gate\_service.py} & Safety gate evaluation before procedure transitions \\
\texttt{approval\_checkpoint.py} & Multi-party approval checkpoints (clinician + system + regulatory) \\
\texttt{site\_verifier.py} & Site capability verification against conformance requirements \\
\bottomrule
\end{tabularx}
\end{table}

The emergency stop controller in \texttt{safety/estop.py} implements a four-state lifecycle (IDLE, TRIGGERED, ACKNOWLEDGED, RECOVERING) with signal propagation to all registered servers. This ensures that any authorized party can halt a robotic procedure immediately, with full state preservation and evidence capture.

\subsection{Test Coverage and Validation}

The repository contains 668 test functions across two test suites:

\begin{table}[H]
\centering
\caption{Test coverage summary across unit and conformance suites.}
\label{tab:test-coverage}
\small
\begin{tabularx}{\textwidth}{@{}L{3cm}R{1.5cm}L{7cm}@{}}
\toprule
\textbf{Category} & \textbf{Tests} & \textbf{Validates} \\
\midrule
Unit (integrations) & 205 & All 34 integration adapters \\
Unit (safety) & 88 & All 8 safety modules \\
Unit (core server) & 32 & 5 MCP server implementations \\
Unit (infrastructure) & 12 & Conformance harness, schema validator \\
Positive conformance & 69 & Correct behavior at Levels 1--3 \\
Adversarial & 57 & AuthZ bypass, chain tampering, PHI leakage, replay attacks \\
Blackbox & 63 & Per-server external conformance \\
Security & 39 & SSRF prevention, token lifecycle, chain integrity \\
Interoperability & 32 & Cross-server trace, schema validation \\
Negative & 26 & Invalid input rejection, unauthorized access \\
Integration & 13 & Multi-server integration scenarios \\
\midrule
\textbf{Total} & \textbf{668} & \\
\bottomrule
\end{tabularx}
\end{table}

\subsection{Integration Adapters}

The repository implements 34 integration adapter modules across six categories in \texttt{integrations/}:

\begin{table}[H]
\centering
\caption{Integration adapter categories and module counts.}
\label{tab:integration-adapters}
\small
\begin{tabularx}{\textwidth}{@{}L{2cm}R{1.2cm}L{8.5cm}@{}}
\toprule
\textbf{Category} & \textbf{Count} & \textbf{Key Adapters} \\
\midrule
FHIR & 9 & HAPI FHIR, SMART on FHIR, de-identification, terminology, patient filter \\
DICOM & 9 & dcm4chee, Orthanc, DICOMweb, RECIST, modality filter \\
Clinical & 3 & Scheduling, e-consent, provenance export \\
Federation & 4 & Coordinator, policy enforcement, secure aggregation \\
Identity & 5 & OIDC, mTLS, KMS, policy engine \\
Privacy & 4 & De-identification pipeline, privacy budget, data residency \\
\midrule
\textbf{Total} & \textbf{34} & \\
\bottomrule
\end{tabularx}
\end{table}

\subsection{Repository Scale}

Table~\ref{tab:repo-scale} summarizes the quantitative scope of the national repository at version 1.2.0:

\begin{table}[H]
\centering
\caption{Repository quantitative summary (v1.1.0).}
\label{tab:repo-scale}
\small
\begin{tabular}{@{}lr|lr@{}}
\toprule
\textbf{Category} & \textbf{Count} & \textbf{Category} & \textbf{Count} \\
\midrule
MCP servers & 5 & Safety modules & 8 \\
Tools (total) & 23 & Benchmark modules & 5 \\
JSON schemas & 13 & Certification tools & 4 \\
Conformance profiles & 8 & Deployment modes & 6 \\
Test functions & 668 & Interop scenarios & 8 \\
Integration adapters & 34 & Actor personas & 6 \\
Specification docs & 9 & Stakeholder guides & 5 \\
ADRs & 7 & Operational docs & 5 \\
Total files & 381 & Total directories & 88 \\
Lines of code (approx.) & 69,800 & Regulatory overlays & 7 \\
\bottomrule
\end{tabular}
\end{table}

\subsection{Deployment Infrastructure}

The system supports six deployment configurations, enabling evaluation at any scale:

\begin{enumerate}[leftmargin=*]
\item \arraybackslash \raggedright \textbf{Individual Docker containers:} Five Dockerfiles in \texttt{deploy/docker/} (e.g., \texttt{Dockerfile.authz}, \texttt{Dockerfile.fhir}).
\item \textbf{All-in-one Docker:} \texttt{Dockerfile.allinone} for single-container evaluation.
\item \textbf{Docker Compose (single-site):} \texttt{deploy/docker-compose.yml} for local deployment.
\item \textbf{Docker Compose (multi-site):} \texttt{deploy/docker-compose.multi-site.yml} for federated testing.
\item \textbf{Kubernetes manifests:} 7 manifests in \texttt{deploy/kubernetes/} for production clusters.
\item \textbf{Helm chart:} \texttt{deploy/helm/trialmcp/} with \texttt{Chart.yaml} and \texttt{values.yaml}.
\end{enumerate} \raggedright

Server configuration is managed through YAML files in \texttt{deploy/config/} (one per server), enabling site-specific customization of storage backends, authentication providers, and privacy settings.

\subsection{SDK and Developer Tooling}

The repository provides dual-language SDKs and developer tools to facilitate adoption:

\begin{itemize}[leftmargin=*]
\item \textbf{Python SDK} (\texttt{sdk/python/trialmcp\_client/}): 8 client modules with 4 middleware components (audit, auth, circuit breaker, retry). Installable via \texttt{pip install -e .} with \texttt{pyproject.toml}.

\item \textbf{TypeScript SDK} (\texttt{sdk/typescript/src/}): 8 client modules with matching middleware stack. Configured via \texttt{package.json} and \texttt{tsconfig.json}.

\item \textbf{CLI} (\texttt{tools/cli/trialmcp\_cli.py}): Commands for \texttt{init}, \texttt{scaffold}, \texttt{validate}, \texttt{certify}, \texttt{schema diff}, and \texttt{config generate}.

\item \textbf{Code generation} (\texttt{tools/codegen/}): Generators for Python models, TypeScript interfaces, and OpenAPI specifications from JSON schemas.
\end{itemize}

Both SDKs include role-scoped example scripts for all six actor types: robot agent, trial coordinator, data monitor, auditor, sponsor, and CRO.

\subsection{Schema and Specification Infrastructure}

The repository defines 13 machine-readable JSON schemas in \texttt{schemas/} using JSON Schema Draft 2020-12. These schemas serve as the normative contract between MCP server implementations and clients, enabling automated validation and code generation:

\begin{table}[H]
\centering
\caption{JSON schemas organized by domain category.}
\label{tab:schemas}
\scriptsize
\begin{tabularx}{\textwidth}{@{}L{4.5cm}L{1.5cm}L{5.5cm}@{}}
\toprule
\textbf{Schema File} & \textbf{Category} & \textbf{Validates} \\
\midrule
\texttt{audit-record.schema.json} & Audit & Audit trail entries \\
\texttt{authz-decision.schema.json} & AuthZ & RBAC evaluation results \\
\texttt{capability-descriptor.schema.json} & Infra & Server capability declarations \\
\texttt{consent-status.schema.json} & Privacy & Patient consent state \\
\texttt{dicom-query.schema.json} & Imaging & DICOM query parameters \\
\texttt{error-response.schema.json} & Infra & Standardized error envelopes \\
\texttt{fhir-read.schema.json} & Clinical & FHIR resource read responses \\
\texttt{fhir-search.schema.json} & Clinical & FHIR search bundle responses \\
\texttt{health-status.schema.json} & Infra & Server health check responses \\
\texttt{provenance-record.schema.json} & Prov & DAG lineage records \\
\texttt{robot-capability-profile} & Robotics & Robot system capabilities \\
\texttt{site-capability-profile} & Infra & Site-level capabilities \\
\texttt{task-order.schema.json} & Robotics & Procedure task orders \\
\bottomrule
\end{tabularx}
\end{table}

Nine normative specification documents in \texttt{spec/} define the complete protocol surface, including \texttt{core.md} (protocol scope, design principles, 5 conformance levels), \texttt{actor-model.md} (6 actors with permission matrices), \texttt{tool-contracts.md} (23 tool signatures with input/output contracts), \texttt{security.md} (RBAC, mTLS, OIDC, SSRF prevention), and \texttt{privacy.md} (HIPAA Safe Harbor, de-identification, consent, data residency).

\subsection{Emergency Stop Implementation}

The emergency stop controller is a critical safety component. The following code excerpt from \texttt{safety/estop.py} shows the four-state lifecycle and signal propagation:

\begin{lstlisting}[caption={Emergency stop lifecycle from \texttt{safety/estop.py}.},label={lst:estop}]
class EStopStatus(Enum):
    IDLE = "IDLE"
    TRIGGERED = "TRIGGERED"
    ACKNOWLEDGED = "ACKNOWLEDGED"
    RECOVERING = "RECOVERING"

class EStopController:
    def trigger_estop(self, procedure_id, reason,
                      triggered_by, current_state=None):
        if self._status == EStopStatus.TRIGGERED:
            raise RuntimeError(
                "E-stop already active")
        signal = EStopSignal(
            signal_id=str(uuid.uuid4()),
            procedure_id=procedure_id,
            reason=reason,
            triggered_by=triggered_by,
            affected_servers=list(
                self._active_servers),
            preserved_state=current_state or {})
        self._active_signal = signal
        self._status = EStopStatus.TRIGGERED
        return signal
\end{lstlisting}

The e-stop lifecycle ensures: (1) only one e-stop can be active at a time, (2) all registered servers receive the stop signal, (3) procedure state is preserved for post-incident analysis, and (4) recovery requires explicit re-authorization. These properties are validated by 88 unit tests in \texttt{tests/test\_safety.py}.

\subsection{Interoperability Testbed}

The interoperability testbed in \texttt{interop-testbed/} provides a multi-site simulation environment with 8 scenarios and 6 actor personas. Each scenario validates a specific cross-site interaction pattern:

\begin{enumerate}[leftmargin=*]
\item \texttt{audit\_replay.py}: Verifies audit trail replay across sites
\item \texttt{cross\_site\_provenance.py}: Validates cross-site provenance chain integrity
\item \texttt{partial\_outage.py}: Tests graceful degradation during partial failures
\item \texttt{robot\_workflow.py}: End-to-end robotic procedure workflow across servers
\item \texttt{schema\_drift.py}: Tests schema version compatibility under evolution
\item \texttt{site\_onboarding.py}: Validates new site onboarding workflow
\item \texttt{state\_overlay.py}: Tests state-level regulatory overlay application
\item \texttt{token\_exchange.py}: Validates cross-site token exchange protocols
\end{enumerate}

The testbed uses Docker Compose (\texttt{interop-testbed/docker-compose.yml}) with mock services for EHR (\texttt{mock\_ehr.py}), PACS (\texttt{mock\_pacs.py}), and identity (\texttt{mock\_identity.py}) to simulate realistic clinical environments.

%% ========================================================================
